% !TeX root = ../../book.tex
\subsection*{可数性}

证明 \Crefrange{cqProveCountableBegin}{cqProveCountableEnd} 中给出的集合是可数的。

\begin{chapex}
\label{cqProveCountableBegin}
\textit{二进}有理数集合 $D = \left\{ \dfrac{a}{2^n} \mid a \in \mathbb{Z}, n \in \mathbb{N} \right\}$
\end{chapex}

\begin{chapex}
所有 $[n] \to \mathbb{Z}$ 函数的集合，其中 $n \in \mathbb{N}$。
\end{chapex}

\begin{chapex}
平方为有理数的所有实数的集合。
\end{chapex}

\begin{chapex}
\label{cqProveCountableEnd}
如下集合：
\[
[(\mathbb{Z} \times \mathbb{Q}) \setminus (\mathbb{N} \times \mathbb{Z})] \cup \{ n \in \mathbb{N} \mid \exists u,v \in \mathbb{N},\, n=5u+6v \} \cup \{ x \in \mathbb{R} \mid x-\sqrt{2} \in \mathbb{Q} \}
\]
\end{chapex}

证明 \Crefrange{cqProveUncountableBegin}{cqProveUncountableEnd} 中给出的集合是不可数的。

\begin{chapex}
\label{cqProveUncountableBegin}
所有 $\mathbb{Z} \to \{ 0,1 \}$ 函数的集合。
\end{chapex}

\begin{chapex}
所有子集 $U \subseteq \mathbb{N}$ 的集合，其中要么 $U$ 是有限的，要么 $\mathbb{N} \setminus U$ 是有限的。
\end{chapex}

\begin{chapex}
\label{cqProveUncountableEnd}
收敛于 $0$ 的所有有理数序列的集合。
\end{chapex}

判断 \Crefrange{cqDetermineIfCountableBegin}{cqDetermineIfCountableEnd} 中给出的集合是可数的还是不可数的，并给出证明。

\begin{chapex}
\label{cqDetermineIfCountableBegin}
所有弱递减函数 $f : \mathbb{N} \to \mathbb{N}$ 的集合 --- 即，对于所有 $m,n \in \mathbb{N}$，如果 $m \le n$，则 $f(m) \ge f(n)$。
\end{chapex}

\begin{chapex}
所有弱递减函数 $f : \mathbb{N} \to \mathbb{Z}$ 的集合。
\end{chapex}

% \begin{definition}
% \label{defPeriodicFunction}
% \index{function!periodic}
% \index{periodic function}
% Given a set $X$ and a subset $U \subseteq \mathbb{R}$ that is closed under addition (that is, $a+b \in U$ for all $a,b \in U$), a function $f : U \to X$ is \textbf{periodic} if there exists some positive $p \in U$ such that $f(x+p)=f(x)$ for all $x \in U$.
% \end{definition}

\begin{chapex}
所有周期函数 $f : \mathbb{Z} \to \mathbb{Q}$ 的集合 --- 即，存在某个整数 $p > 0$，对于所有 $x \in \mathbb{Z}$, $f(x+p) = f(x)$。
\end{chapex}

\begin{chapex}
所有周期函数 $f : \mathbb{Q} \to \mathbb{Z}$ 的集合 --- 即，存在某个有理数 $p > 0$，对于所有 $x \in \mathbb{Q}$, $f(x+p) = f(x)$。
\end{chapex}

\begin{chapex}
\label{cqDetermineIfCountableEnd}
对于有理数 $a_0,a_1,\dots,a_d$，其中 $a_d \ne 0$，满足 $a_0 + a_1x + \cdots + a_dx^d = 0$ 的所有实数 $x$ 的集合。
\end{chapex}

\begin{chapex}
对于所有 $a \in \mathbb{R}$ 且 $\varepsilon > 0$，如果 $(a-\varepsilon, a+\varepsilon) \cap U$ 非空，则说子集 $D \subseteq \mathbb{R}$ 是\textit{稠密的} --- 直观上，这意味着 $U$ 中存在与任何实数任意接近的元素。$\mathbb{R}$ 的稠密子集一定是不可数的吗？
\end{chapex}

\subsection*{基数}

判断 \Crefrange{cqFindCardinalityBegin}{cqFindCardinalityEnd} 给出集合的基数等于 $\aleph_0$，还是等于 $\mathfrak{c}$，抑或是大于 $\mathfrak{c}$。

\begin{chapex}
\label{cqFindCardinalityBegin}
$\mathcal{P}(\mathbb{R})$。
\end{chapex}

\begin{chapex}
所有 $\mathbb{R}$ 的有限子集的集合。
\end{chapex}

\begin{chapex}
所有 $\mathbb{R}$ 的可数无限子集的集合。
\end{chapex}

\begin{chapex}
$\mathbb{R} \times \mathbb{R}$。
\end{chapex}

\begin{chapex}
所有函数 $f : \mathbb{N} \to \mathbb{R}$ 
The set of all functions $f : \mathbb{N} \to \mathbb{R}$ such that $f(n) = 0$ for all but finitely many $n \in \mathbb{N}$.
\end{chapex}

\begin{chapex}
所有 $\mathbb{Q} \to \mathbb{R}$ 函数的集合。
\end{chapex}

\begin{chapex}
\label{cqFindCardinalityEnd}
所有 $\mathbb{R} \to \mathbb{R}$ 函数的集合。
\end{chapex}

\begin{chapex}
证明 $\mathbb{R}^2$ 中存在一个两两不相交的圆的集合 $\mathcal{C}$，使得 $|\mathcal{C}| = \mathfrak{c}$；形式上，圆是以下形式的子集 $C \subseteq \mathbb{R}^2$
\[ C = \{ (x,y) \mid (x-a)^2 + (y-b)^2 = r^2 \} \]
其中 $a,b \in \mathbb{R}$ 且 $r > 0$。
\end{chapex}

\begin{chapex}
证明 $\mathbb{R}^2$ 中不存在一个两两不相交的圆盘的集合 $\mathcal{D}$，使得 $|\mathcal{D}| = \mathfrak{c}$；形式上，圆盘一下形式的子集 $D \subseteq \mathbb{R}^2$
\[ D = \{ (x,y) \mid (x-a)^2 + (y-b)^2 \le r^2 \} \]
其中 $a,b \in \mathbb{R}$ 且 $r > 0$。
\hintlabel{cqDiscsInPlane}{%
首先证明任意圆盘 $D \subseteq \mathbb{R}^2$ 包含一个坐标都是有理数的点。
}
\end{chapex}

\subsection*{基数算术}

\begin{chapex}
证明对于所有基数 $\kappa$, $\kappa + 0 = \kappa \cdot 1 = \kappa^1 = \kappa$。
\end{chapex}

\begin{chapex}
证明对于所有基数 $\kappa$, $\kappa^0 = 1^{\kappa} = 1$。
\end{chapex}

\begin{chapex}
设 $\kappa$ 为无限基数，且满足 $\kappa \cdot \kappa = \kappa$。证明 $\kappa^{\kappa} = 2^{\kappa}$。
\end{chapex}

\begin{chapex}
证明 $\mathfrak{c}^{\mathfrak{c}} = 2^{2^{\aleph_0}}$。
\end{chapex}

\begin{chapex}
设 $\kappa$, $\lambda$ 和 $\mu$ 为基数。
\begin{enumerate}[(a)]
\item 证明如果 $\kappa \le \lambda$，则 $\kappa + \mu \le \lambda + \mu$。
\item 假设 $\kappa + \mu \le \lambda + \mu$。一定有 $\kappa \le \lambda$ 吗？
\end{enumerate}
\end{chapex}

\begin{chapex}
设 $\kappa$, $\lambda$ 和 $\mu$ 为基数。
\begin{enumerate}[(a)]
\item 证明如果 $\kappa \le \lambda$，则 $\kappa \cdot \mu \le \lambda \cdot \mu$。
\item 假设 $\kappa \cdot \mu \le \lambda \cdot \mu$ 且 $\mu > 0$。一定有 $\kappa \le \lambda$ 吗？
\end{enumerate}
\end{chapex}

\begin{chapex}
设 $\kappa$, $\lambda$ 和 $\mu$ 为基数。
\begin{enumerate}[(a)]
\item 证明如果 $\kappa \le \lambda$，则 $\kappa^{\mu} \le \lambda^{\mu}$。
\item 假设 $\kappa^{\mu} \le \lambda^{\mu}$ 且 $\mu > 0$。一定有 $\kappa \le \lambda$ 吗？
\end{enumerate}
\end{chapex}

\begin{chapex}
设 $\kappa$, $\lambda$ 且 $\mu$ 为基数。
\begin{enumerate}[(a)]
\item 证明如果 $\kappa \le \lambda$，则 $\mu^{\kappa} \le \mu^{\lambda}$。
\item 假设 $\mu^{\kappa} \le \mu^{\lambda}$ 且 $\mu > 1$。一定有 $\kappa \le \lambda$ 吗？
\end{enumerate}
\end{chapex}

\begin{definition}
给定基数 $\kappa$ 和 $\lambda$，定义\textbf{二项式系数} $\dbinom{\kappa}{\lambda}$ 如下
\[ \dbinom{\kappa}{\lambda} = |\{ U \subseteq [\kappa] \mid |U| = \lambda \}| \]
\end{definition}

\begin{chapex}
设 $\kappa$ 为基数。证明
\[ \dbinom{\aleph_0}{\kappa} = \begin{cases} 1 & \text{如果 } \kappa = 0 \\ \aleph_0 & \text{如果 } \kappa \in \mathbb{N} \text{ 且 } \kappa > 0 \\ 2^{\aleph_0} & \text{如果 } \kappa = \aleph_0 \\ 0 & \text{如果 } \kappa \not\in \mathbb{N} \cup \{ \aleph_0 \} \end{cases} \]
\end{chapex}

\begin{chapex}
对于 $\kappa \in \mathbb{N} \cup \{ \aleph_0, \mathfrak{c} \}$，找到 $\dbinom{\mathfrak{c}}{\kappa}$ 的值。
\end{chapex}

\begin{chapex}
对于所有基数 $\kappa$，证明 $\displaystyle \sum_{\lambda \le \kappa} \dbinom{\kappa}{\lambda} = 2^{\kappa}$。
\end{chapex}

\begin{chapex}
通过推广自然数 $n$ 的阶乘 $n!$，定义基数 $\kappa$ 的阶乘 $\kappa!$。找到 $\aleph_0!$ 和 $\mathfrak{c}!$ 的表达式。
\hintlabel{cqFactorialOfCardinalNumber}{%
组合地思考该问题。不应该尝试泛化阶乘的递归定义（\Cref{defFactorialRecursive}）；而且\textit{更不该}尝试泛化阶乘的非正式定义`$n! = 1 \times 2 \times \cdots \times n$'。
}
\end{chapex}

\subsection*{判断题}

\tfquestiontext{cqInfinityTFBegin}{cqInfinityTFEnd}

\begin{chapex} % True
\label{cqInfinityTFBegin}
对于任意两个可数无限集 $X$ 和 $Y$，存在双射 $X \to Y$。
\end{chapex}

\begin{chapex} % False
对于任意两个不可数无限集 $X$ 和 $Y$，存在双射 $X \to Y$。
\end{chapex}

\begin{chapex} % False
可数个可数集的积是可数的。
\end{chapex}

\begin{chapex} % True
每一个可数无限集可以分割成无限多个无限子集。
\end{chapex}

\begin{chapex} % False
$\mathbb{N}$ 有可数多个无限子集
\end{chapex}

\begin{chapex} % True
集合 $\mathbb{N} \times \mathbb{Z} \times \mathbb{Q}$ 是可数无限集。
\end{chapex}

\begin{chapex} % False
$\aleph_0$ 是最小的基数。
\end{chapex}

\begin{chapex} % False
存在不同基数的集合 $X$ 和 $Y$，它们具有相同数量的有限子集。
\end{chapex}

\begin{chapex} % True
存在不同基数的集合 $X$ 和 $Y$，它们具有相同数量的可数无限子集。
\end{chapex}

\begin{chapex} % False
对于所有集合 $X$，如果 $U \subsetneqq X$，则 $|U| < |X|$。
\end{chapex}

\begin{chapex} % False
\label{cqInfinityTFEnd}
对于所有基数 $\kappa > 0$, $2^{\kappa} < 3^{\kappa}$。
\end{chapex}

\subsection*{可能性问题}

\asnquestiontext{cqInfinityASNBegin}{cqInfinityASNEnd}

\begin{chapex} % Always
\label{cqInfinityASNBegin}
设 $X$ 为集合，并假设存在单射 $\mathbb{N} \to X$。则存在单射 $\mathbb{Q} \to X$。
\end{chapex}

\begin{chapex} % Never
设 $X$ 为集合。则 $\mathcal{P}(X)$ 为可数无限集。
\end{chapex}

\begin{chapex} % Always
设 $X$ 为集合，并设 $\kappa$ 和 $\lambda$ 为基数，假设 $\lambda \le \kappa = |X|$。则 $X$ 有大小为 $\lambda$ 的子集。
\end{chapex}

\begin{chapex} % Always
设 $X$ 和 $Y$ 为集合，其中 $|X| \le |Y|$。则 $|\mathcal{P}(X)| \le |\mathcal{P}(Y)|$。
\end{chapex}

\begin{chapex} % Never
设 $n \in \mathbb{N}$ 其中 $n \ge 2$。则 $\aleph_0^n > \aleph_0$。
\end{chapex}

\begin{chapex} % Always
\label{cqInfinityASNEnd}
设 $\kappa$ 为基数，$\kappa \ge \aleph_0$，则 $\aleph_0^{\kappa} > \aleph_0$。
\end{chapex}

\subsection*{刁钻问题}

\begin{chapex}
存在基数 $\kappa$ 使得 $\aleph_0 < \kappa < 2^{\aleph_0}$ 吗？
\hintlabel{cqContinuumHypothesis}{%
不要浪费太多时间在这个问题上 --- 它已被证明独立于通常的数学公理（参见 \Cref{secZFC}），这意味着答案为`真'和答案为`假'都无法被证明（因此这两个答案都不会导致矛盾）。答案为`假'称为\textit{连续统假设}。
}
\end{chapex}
